% !TeX root = ../main.tex
%------------------------------------------------------------------------------
% Digital Engineering and the Lifecycle Digital Twin
%
% Placed in Part I: the information backbone is established before the ground
% chapters, because every later chapter reads from and writes to it.
%------------------------------------------------------------------------------
\chapter{Digital Engineering and the Lifecycle Digital Twin}
\label{ch:digital-twin}

Information management has become a defining feature of underground practice,
and the object at its centre is a combined representation of structure and
ground rather than a geometric model of the structure alone. Such a model carries strata
and inferred boundaries, hydrogeological units, design parameters and their
uncertainty, construction sequence and support timing, quantities and carbon
factors, and monitoring and inspection records. Its purpose is to allow
design, construction, carbon accounting, risk management and operation to draw
on a single description of the project rather than on separate and divergent
ones.

This chapter comes before the ground chapters because the
characterisation, design, construction and operation chapters that follow all
read from and write to this model. Keeping one current
model is what makes the lifecycle argument of Chapter~\ref{ch:challenges}
workable. The subject here is information flow and
decision traceability, not software.

\section{The geo-integrated asset model}
\label{sec:geo-integrated-model}

\subsection{Why geometry alone is insufficient}
\label{sec:beyond-geometry}

Building information modelling arrived in underground engineering through
architecture and building services, and it carried their assumptions with it.
A building model describes objects that will be manufactured to specification:
their geometry is known, their properties are declared, and the uncertainty
lies in construction rather than in the medium. None of that holds
underground. The dominant object is the ground, its geometry is interpolated
between sparse observations, its properties are inferred through correlation
chains, and its uncertainty is irreducible.

A model that represents only the structure therefore omits the governing
variable, and it encourages the ground to be treated as the space
around the tunnel rather than as the material being designed in. The
distinguishing feature of a useful underground model is that geology,
hydrogeology and their uncertainty are modelled objects with the same standing
as the lining.

The second departure is linearity. Underground assets are alignments, and
their information is naturally indexed by chainage rather than by storey or
grid line. Ground class, support class, monitoring section, quantities and
as-built records are all functions of chainage (Figure~\ref{fig:dt-chainage}), and a model that cannot
express that relationship forces every downstream user to rebuild it.

\begin{figure}[htbp]
\centering
\datafigure[\textwidth]{ch02/dt-chainage}
\caption{One model indexed by chainage (illustrative): geological unit, with a
fault zone in carmine; ground class; support class, extended beyond each change
in ground class; monitoring sections, densified through the fault; and lining
carbon per metre relative to the lightest support. A change in any row is
traceable to the others through the common coordinate.}
\label{fig:dt-chainage}
\end{figure}

\subsection{What the model must contain}
\label{sec:model-contents}

A model that serves the whole lifecycle integrates, in one structure:

\begin{itemize}
 \item infrastructure function and system interfaces: what service the
 asset provides and what it connects to;
 \item alignment, geometry, lining, support, shafts, stations and fit-out;
 \item borehole data, geological interpretation, hydrogeological units,
 design parameters and their uncertainty;
 \item construction method, sequence, support timing, machine performance and
 haulage assumptions;
 \item quantities, carbon factors, energy factors and the declared system
 boundary;
 \item monitoring, inspection, condition, damage and maintenance records.
\end{itemize}

The first three establish what is being built and in what; the fourth records
how; the fifth and sixth are what make carbon assessment and asset management
possible without a separate data-gathering exercise. Omitting any one of them
breaks a link in the chain that
carries a decision back to its evidence.

\subsection{Level of detail as a decision variable}
\label{sec:level-of-detail}

Detail should follow the decision. Early planning needs system-level
quantities and formation-scale geometry; detailed design needs component
dimensions and material properties; construction control needs metre-scale
ground classification; operation needs asset registers and condition records.
These are different models of the same asset, and the common error is to build
all of them at the finest available resolution.

That error has two costs: directly, in modelling effort
that no decision requires, and indirectly, because a model detailed beyond its
evidence base implies a precision the data cannot support. A ground surface
drawn to millimetre tolerance between boreholes 200\,m apart describes
the interpolator rather than the ground (Figure~\ref{fig:dt-interpolation}).

\begin{figure}[htbp]
\centering
\datafigure[\textwidth]{ch02/dt-interpolation}
\caption{Depth to rock between boreholes 200~m apart (synthetic). A smooth and
a linear interpolant both honour the boreholes exactly and agree closely, yet
both miss the trough near 300~m. The 90\,\% band of the smooth interpolant
collapses at each borehole and widens between them; even it excludes the
trough, which no interpolation of these five holes can recover.}
\label{fig:dt-interpolation}
\end{figure}

The practical discipline is to state, for each model, the decisions it exists
to support and the tolerance those decisions require, and to add no detail
that serves neither.

\section{Standards, environments and the analysis loop}
\label{sec:standards-and-loop}

\subsection{Interoperability and data standards}
\label{sec:data-standards}

An underground asset outlives the software used to design it, usually by
several decades and often by a century. Information held only in a proprietary
format is therefore liable to become unreadable within the life of the asset,
and the cost falls on the operator rather than on the designer who chose the format.

Open exchange schemas exist to prevent this. The Industry Foundation Classes
(IFC) family provides the established route for building and infrastructure
data, extended for linear underground assets by tunnel-specific schema
developments, with
alignment geometry and ground objects represented alongside structure.
Information management standards in the ISO 19650 series govern how that
information is named, versioned, approved and handed over~\citep{iso19650}.

Neither standard makes a model correct. They make it legible to
someone who did not build it, which is the condition for the model surviving
the transition from designer to contractor to operator.

\subsection{Common data environments}
\label{sec:cde}

A common data environment is the agreed single location for project
information, with defined states (work in progress, shared, published,
archived) and an audit trail of who changed what and when. Its engineering
significance is that it establishes which version of the ground model a
calculation was based on, and therefore whether a later divergence between
prediction and observation reflects a change in the ground, a change in the
model, or an error.

The contractual significance is equal. Where model information forms part of
the contract, the status of a given version determines whether a contractor
was entitled to rely on it. Ambiguity about that status later becomes the
subject of a claim.

\subsection{Forward prediction and backward updating}
\label{sec:forward-backward}

The loop that gives the model its value has two directions
(Figure~\ref{fig:dt-update-loop}). Forward, analysis
models are instantiated from the asset and ground models (settlement,
support interaction, lining response, excavation sequence) rather than
having their geometry rebuilt by hand for each study. Backward, observations
from construction and monitoring update parameters, refine the ground
interpretation and revise the design ahead of the face. The formal mechanism
by which an observation revises a parameter is set out in
Section~\ref{sec:bayesian-updating}.

\begin{figure}[htbp]
\centering
\datafigure[\textwidth]{ch02/dt-update-loop}
\caption{Forward prediction and backward updating. Analysis models are
instantiated from the asset and ground model, and their predictions are
compared with observation. An observation updates the model only after the
movement already present before construction has been removed and the
remainder attributed to the works (Figure~\ref{fig:movement-attribution}).}
\label{fig:dt-update-loop}
\end{figure}

Both directions are necessary. Forward analysis alone produces a prediction
nobody checks; backward updating alone produces a record nobody uses. The loop
must also stay active after construction, because the operational phase
generates the longest and most complete record of how the asset
behaves.

Keeping the loop active is an organisational problem more than a technical one. It
requires that someone owns the model, that updating it is a scheduled activity
rather than a voluntary one, and that the design team remains engaged long
enough to receive the evidence that tests its assumptions.

Backward updating begins before construction. An observation is not by
itself evidence about the excavation, because ground moves for reasons
unrelated to the works (groundwater abstraction, consolidation
under existing load, seasonal shrink and swell). A reading that is not
referred to the movement already under way attributes all of that movement
to the tunnel. Establishing the pre-construction field is therefore part of the
loop rather than preparation for it, and
Figure~\ref{fig:movement-attribution} shows the separation it makes possible.

\begin{figure}[htbp]
\centering
\datafigure{ch02/twin-anomaly}
\caption[Separating construction movement from movement already present]{%
Attribution as the first step of backward updating. Each mark is a building
above an illustrative twin-tunnel corridor, placed by its cumulative satellite
radar movement over the seven years before construction against its cumulative
movement at the end of the first year of construction, and shaded by the
construction contribution. A building on the 1:1 line has not moved during
construction. Synthetic demonstrator data (Appendix~\ref{app:software}).}
\label{fig:movement-attribution}
\end{figure}

\section{Digital twins}
\label{sec:digital-twins}

\subsection{Four levels of maturity}
\label{sec:twin-maturity}

The term \emph{digital twin} is used loosely enough to be unhelpful without
qualification. The useful distinction between a model and a twin is the
feedback loop, and four levels can be separated by what each is able to
answer:

\begin{description}[leftmargin=28mm,style=nextline]
 \item[Descriptive] \emph{What is the asset?} Geometry, geology, systems,
 materials and historical records organised into a structured asset model.
 This level is a record; it does not change unless someone changes it.
 \item[Reflective] \emph{What is happening now?} Monitoring, machine data,
 inspection and condition observations are ingested to validate the model
 and diagnose the current state.
 \item[Predictive] \emph{What will happen next?} Physical, statistical,
 machine-learning and surrogate models (Section~\ref{sec:surrogate-models})
 forecast settlement, convergence,
 leakage, deterioration, damage, energy, carbon and functionality.
 \item[Prescriptive] \emph{What should be done?} Intervention options are
 compared on risk, recovery, cost, time, carbon and service outcome, and a
 recommendation is produced.
\end{description}

Each level requires the one before it (Figure~\ref{fig:dt-maturity}). The step to reflective adds a class of
data the descriptive level did not need; the steps to predictive and
prescriptive add models and decision criteria to data already being
collected. A predictive model built without a reflective layer is a
simulation rather than a twin, because it forecasts from assumptions and not
from the asset's measured state.

\begin{figure}[htbp]
\centering
\datafigure[\textwidth]{ch02/dt-maturity}
\caption{The four levels of twin maturity, each defined by one question and
each requiring the level before it. The step to reflective adds measured
data; the steps to predictive and prescriptive add forecasting models and
decision criteria to data already being collected.}
\label{fig:dt-maturity}
\end{figure}

\subsection{Analytics and machine learning}
\label{sec:analytics}

Mechanised tunnelling generates a continuous, high-resolution record of the
machine's interaction with the ground (thrust, torque, penetration rate,
cutterhead rotation, face pressure, grout volumes), logged ring by ring. That
record is a geotechnical measurement as well as an operational one, and it is
usually the densest ground data a project ever collects.

Analytics on that record supports three distinct tasks: inferring ground
conditions behind the face, detecting departure from expected behaviour in
real time, and predicting performance ahead. Supervised learning has been
applied to all three, along with automated photogrammetric and LiDAR face
mapping and lining defect detection from imagery.

Two cautions govern their use. The first is transferability: a model trained
on one machine in one geological setting encodes both, and its performance on
a different drive has to be tested (Section~\ref{sec:ml-validation}). The
second is extrapolation: a prediction outside the training envelope (the fault
zone, the mixed face, the unexpected boulder) is returned with the same
apparent confidence as one inside it (Section~\ref{sec:ml-drift}).

\subsection{Governance and professional responsibility}
\label{sec:data-governance}

When a model informs a safety-critical decision, three questions must have
answers: where did this value come from, who approved it, and what would have
to be true for it to be wrong. Provenance, approval status and stated
assumptions are therefore part of the model, not metadata about it.

This bears directly on emerging methods. Digital twins, machine-learning
prediction and prescriptive maintenance are presented in this book as
forward-looking capabilities, and they should be presented that way in design
documentation until supported by mature project evidence. Separating
established practice from developing method
allows a reviewer to know which parts of a submission rest on a case base and
which rest on a model.

Accountability does not transfer to the model. An engineer who cannot explain
why a prediction is reasonable, independently of the tool that produced it,
has not yet completed the design.

\section{Assurance of data-driven methods}
\label{sec:ml-assurance}

Section~\ref{sec:data-governance} states where accountability sits; this
section states what has to be done to discharge it. A statistical model informing a
support class needs an assurance regime as specific as the one
Chapter~\ref{ch:numerical-modelling} sets out for a numerical model, but it
cannot use the same one: verification against an analytical solution, mesh
convergence and equilibrium checks have no counterpart here. They are replaced
by a discipline about data, which is unfamiliar enough in geotechnical
practice to be set out explicitly.

\subsection{What validation means}
\label{sec:ml-validation}

Validation means demonstrating performance on data the model has not seen.
The difficulty lies in what counts as unseen, and geotechnical data
makes that harder than it is in most fields.

Samples from one borehole are not independent of one another. A model trained
on part of a hole and tested on the rest is tested on material it has
effectively already seen. It will report an accuracy that it does not achieve on a
new hole, which is the only situation in which it will be used. The
unit of splitting must therefore be the unit of future application: whole
boreholes for a lithology classifier, whole drives or machines for a model
built on tunnelling machine data, whole projects where the claim is that the
model transfers. Splitting at the sample is a recurrent defect in
published geotechnical machine learning, and it inflates reported accuracy
substantially. Figure~\ref{fig:validation-split} shows the split made at the
machine, on records where the answer for the withheld machine is known.

\begin{figure}[htbp]
\centering
\datafigure{ch02/twin-transfer}
\caption[Validation on a split at the unit of future application]{%
Validation on a split at the unit of future application. A settlement model
trained on seven tunnelling machines along an illustrative corridor is tested
on the eighth after each training round (solid). The dashed curve is the same
model trained only on the withheld machine's limited records. The dotted line
is the transferred model after calibration on those records, 3.28~mm, a single
final result. Beyond the first few rounds the transferred model has the lower
error. Synthetic demonstrator data (Appendix~\ref{app:software}).}
\label{fig:validation-split}
\end{figure}

The metric must match the decision. Overall accuracy is a poor choice where
classes are imbalanced, and geotechnical classes always are: the fault gouge
that governs the design is the rarest label in the training set, and a model
that never predicts it can still be right nineteen times in twenty. The
costs are also asymmetric (failing to predict weak ground is not the same error
as predicting it where it is absent), so a metric that treats the two alike
does not measure what matters to the project. Recall on a class, the
proportion of the cases truly belonging to it that the model finds, reported
for the consequential classes alongside the full table of predicted against
actual class (the confusion matrix), states something a design reviewer can
act on.

Finally, the domain of validity must be stated as deliberately as it is for
the closed-form solutions of Appendix~\ref{app:closed-form}. A model carries
the geology, the equipment and the practices of the data it was trained on.
Its performance elsewhere is unknown until it is tested,
and a submission should record the training envelope so
that a reviewer can see when a prediction lies outside it.

\subsection{Drift and the operating model}
\label{sec:ml-drift}

A model embedded in a twin cannot be checked once and then left. Its
inputs shift: the machine wears, the operators change, the ground changes, the
instruments are replaced or recalibrated. Performance decays, and
because the model continues to return confident predictions, the decay
is not seen unless it is measured.

Monitoring the model is therefore a separate activity from monitoring the
asset. Two things are watched. The first is the input distribution: when
current inputs no longer resemble the training data, predictions are
extrapolations whatever their apparent confidence, and that comparison can be
made continuously without waiting for ground truth. The second is realised
error, wherever an observation eventually reveals what the prediction should
have been: the face that is mapped, the core that is logged, the
convergence that is measured. Predictions should be logged with their
timestamps so this comparison remains possible; a prediction not recorded
cannot be scored later.

Retraining then needs a trigger and a procedure, both agreed in advance:
what degradation prompts it, who authorises it, and how the retrained model is
validated before it replaces the one in service. Unrecorded retraining is the
worst arrangement, because the model that produced a past decision
can no longer be reconstructed, and a decision whose basis cannot be
reconstructed cannot be defended.

\subsection{What a submission must contain}
\label{sec:ml-submission}

A reviewer asked to accept a data-driven result needs the same things they
need for any other analysis: what was assumed, what was done, and what would
have to be true for it to be wrong. In this context that resolves to six
items, and their absence should be treated as it would be for a numerical
analysis submitted without its boundary conditions.

The six items are: the training data and its provenance, including which projects,
machines and formations it came from; the quantity predicted, distinguished from
the quantity of interest where the two differ (a model predicting logged
weathering grade is not predicting strength); the validation protocol, stating
the splitting unit; the performance on consequential classes rather than
overall; the domain of validity and how the model behaves outside it; and the
named individual who accepts the result, who cannot be the model or
its supplier.

None of this makes a statistical prediction equivalent to a measurement.
It makes the prediction reviewable, which is the
condition for using it in a design. The accountability rule of
Section~\ref{sec:data-governance} still governs.

\section{Generative methods and their present scope}
\label{sec:generative-methods}

Generative language models have entered general professional use faster than any
other tool this book describes, and they change quickly enough that detailed
comment would soon be out of date.
The position taken here is therefore narrow, and follows the rule
stated in Section~\ref{sec:data-governance}: emerging methods are presented as
developing capability until a project evidence base exists.

The applications with a defensible basis today are on the document side rather
than the engineering side, and they are still of practical value. Underground projects
generate a very large unstructured record over decades (factual reports, design
statements, shift logs, instrumentation reviews, correspondence),
and the practical difficulty has always been that the answer to a
present question usually exists somewhere in that record and cannot be found.
Retrieval across a corpus, summarisation of what a document says
about a specific reach, and consistency checking of a submission against a
specification are tasks of that kind. They are useful because the
output is checkable: the source document either says what was reported or it
does not, and the engineer verifies it in the ordinary way.

The claims on the engineering side are less well supported. Predicting ground
behaviour, selecting support, or generating design from a description are
proposed regularly and, at the time of writing, rest on demonstration rather
than on a record of use in delivered projects. Two properties argue for
caution beyond the general reservation about new methods. These models produce
fluent output whether or not they have a basis for it, so the usual signal
that something is wrong (output that looks wrong) is absent. They also do
not reliably report their own uncertainty, which removes the property that
makes the statistical methods of Section~\ref{sec:ml-assurance}
governable: a validated classifier states how often it is wrong and on which classes,
and that is the basis on which a reviewer accepts it.

The governing requirement is therefore unchanged. Any output used in a design must be traceable to a
source that can be checked, attributable to a person who accepts it, and
reproducible well enough that the basis of a past decision can be
reconstructed. A generative tool that assists an engineer in finding and
organising evidence satisfies those conditions. One that supplies
a conclusion the engineer cannot independently justify does not
(Section~\ref{sec:data-governance}), and improvement in the tool does not
change that.

\chapterreferences
